\documentclass[12pt]{article}
\usepackage{xr}
\externaldocument{"../Book_II/chapter_6"}
\externaldocument{"../Book_II/chapter_8"}
\externaldocument{chapter_1}
\externaldocument{chapter_2}
\usepackage{algorithm}
\usepackage{algpseudocode}
\algnewcommand\algorithmicforeach{\textbf{for each}}
\algdef{S}[FOR]{ForEach}[1]{\algorithmicforeach\ #1\ \algorithmicdo}
\usepackage[a4paper, margin=1in]{geometry}
\usepackage{tikz}
\usetikzlibrary{positioning, arrows.meta, decorations.pathreplacing, calc, backgrounds}
\usepackage{amsmath,amssymb}
\usepackage{hyperref}
\usepackage{amsthm}
\usepackage{pdflscape} % For the landscape environment

\usepackage[margin=1in]{geometry}
\usepackage{amsmath,amssymb,amsthm,amsfonts}
\usepackage{graphicx}
\usepackage{hyperref}
\usepackage{enumitem}
\setlist[enumerate]{align=left,labelindent=0pt}
\hypersetup{colorlinks=true, linkcolor=blue, urlcolor=blue, citecolor=blue}
% Chapter-specific theorem environments
\theoremstyle{definition}
\newtheorem{definition}{Definition}[section]
\newtheorem{remark}{Remark}[section]
\theoremstyle{plain}
\newtheorem{theorem}{Theorem}[section]
\newtheorem{axiom}{Axiom}[section]
\newtheorem{lemma}{Lemma}[section]
\newtheorem{proposition}{Proposition}[section]
\newtheorem{corollary}{Corollary}[section]

\title{\textbf{Appendix: Details of ATRT proof}}
\author{Dmytro Surdu}
\begin{document}
\maketitle

\section{Adaptive Traceability Resolution Theorem: Algorithm and Lemmas}

\subsection{Promise Verification via Effective Measurability Lemma}

\begin{lemma}[Promise Verification via Effective Measurability]
\label{lem:promise-verification}
Let \(\mathfrak T\) satisfy (F), (S), and (C) and (EM1)--(EM3). Fix \(z\in\mathcal Z^{\mathrm{in}}\), threshold \(\varepsilon>0\), and candidate gap \(\gamma>0\). Define
\[
f(\Theta):=\operatorname{dist}(U(\Theta,z),\mathcal B),\qquad
f^*:=\min_{\Theta\in\mathcal M} f(\Theta),
\]
and let \(\omega(\cdot)\) be the continuity modulus from Lemma~\ref{lem:modulus-bound}. Then there exists an explicit adaptive procedure that maintains computable lower and upper bounds \(L\le f^*\le U\) such that:
\begin{enumerate}
  \item \(U\) is obtained by sampling candidate states and taking \(U := \min_{\Theta_R} f(\Theta_R)\) over representatives \(\Theta_R\);
  \item \(L\) is obtained by region exclusion via continuity: for each region \(R\) with representative \(\Theta_R\), define \(L_R := f(\Theta_R) - \omega(\operatorname{diam}(R))\), and set \(L := \min_R L_R\);
  \item The procedure terminates in finite time once one of the following holds:
    \begin{itemize}
      \item \(U\le\varepsilon\), certifying \(f^*\le\varepsilon\);
      \item \(L\ge\varepsilon+\gamma\), certifying \(f^*\ge\varepsilon+\gamma\);
      \item \(U-L<\gamma\), in which case \(f^*\notin (\varepsilon,\varepsilon+\gamma)\) and the dichotomy ``either \(f^*\le\varepsilon\) or \(f^*\ge\varepsilon+\gamma\)'' is numerically validated.
    \end{itemize}
\end{enumerate}
In particular, whenever the promised gap condition holds, the procedure reaches one of the stopping criteria in finite time and thus does not rely on assuming the dichotomy a priori without evidence.
\end{lemma}

\begin{proof}
By (C) the space \(\mathcal M\) is compact and \(U\) is continuous, so \(f(\Theta)=\operatorname{dist}(U(\Theta,z),\mathcal B)\) is continuous and attains its minimum \(f^*\). We initialize a finite cover \(\mathcal R_0\) of \(\mathcal M\) by compact regions (for example, axis-aligned boxes) and associate to each active region \(R\) a representative point \(\Theta_R\in R\). Using (EM1)--(EM3) and the known modulus \(\omega\), we carry out the following at each refinement iteration:

\textbf{Upper bound \(U\):} For each representative \(\Theta_R\) compute \(f(\Theta_R)\) to sufficient precision; since \(f^*\le f(\Theta_R)\), define \(U := \min_{R} f(\Theta_R)\). This is a valid upper bound.

\textbf{Lower bound \(L\):} For each region \(R\), Lemma~\ref{cor:region-lower} gives that
\[
\min_{\Theta\in R} f(\Theta) \ge f(\Theta_R) - \omega(\operatorname{diam}(R)) =: L_R.
\]
Set \(L := \min_R L_R\); then \(L\le f^*\). Thus \(L\) is a valid lower bound.

\textbf{Refinement:} For any region \(R\) that is not yet resolved (i.e., neither excluded by \(L_R>\varepsilon+\gamma\), nor containing a witness with \(f(\Theta_R)\le \varepsilon\)) subdivide it into smaller-diameter subregions, update their representatives, and recompute the corresponding \(f\)-values and local lower bounds. Because \(\omega(\operatorname{diam}(R))\to 0\) as \(\operatorname{diam}(R)\to 0\), the gap \(U-L\) strictly decreases under sufficient refinement unless already one of the stopping conditions is met.

\textbf{Termination conditions:}
\begin{itemize}
  \item If \(U\le \varepsilon\), then there exists a sample \(\Theta_R\) with \(f(\Theta_R)\le\varepsilon\), so \(f^*\le\varepsilon\).
  \item If \(L\ge \varepsilon+\gamma\), then for all \(\Theta\in\mathcal M\), \(f(\Theta)\ge f^*\ge L\ge \varepsilon+\gamma\), so \(f^*\ge \varepsilon+\gamma\).
  \item If \(U-L<\gamma\) while neither of the above holds, then \(f^*\) cannot lie in the open interval \((\varepsilon,\varepsilon+\gamma)\), because that would force both \(U>\varepsilon\) and \(L<\varepsilon+\gamma\) with gap at least \(\gamma\); the current enclosures \(L\le f^*\le U\) contradict \(f^*\in(\varepsilon,\varepsilon+\gamma)\). Hence the promise ``either \(f^*\le\varepsilon\) or \(f^*\ge\varepsilon+\gamma\)'' is validated numerically, and a definitive decision follows by testing whether \(U\le\varepsilon\) or \(L\ge\varepsilon+\gamma\).
\end{itemize}

Because the bounds \(U\) and \(L\) converge to \(f^*\) as regions are refined (compactness + continuity + shrinking modulus), one of the termination conditions is reached in finite time whenever the dichotomy holds. Thus the procedure achieves the claimed promise verification without assuming it silently. \(\blacksquare\)
\end{proof}

\begin{corollary}[Gating ATRT on Verified Promise]
\label{cor:promise-gate}
Before executing the core dichotomous branch of ATRT (deciding traceable vs non-traceable with gap), run the procedure of Lemma~\ref{lem:promise-verification}. Proceed with the corresponding branch only once one of the certified conditions \(U\le\varepsilon\), \(L\ge\varepsilon+\gamma\), or \(U-L<\gamma\) is observed. This guarantees the gap assumption is either confirmed or refuted explicitly and is not used implicitly.
\end{corollary}


\subsection{SMT Encoding Templates for Region-Level Traceability Checks}

Let \(z\in\mathcal Z^{\mathrm{in}}\) be the input under test, and define 
\[
f(\Theta)\;:=\;\operatorname{dist}(U(\Theta,z),\mathcal B)
\]
with all computations interpreted to finite precision. Let \(R\subseteq\mathcal M\) be a compact region (e.g., a box or cell) represented by constraints \(C_R(\Theta)\).

\paragraph{Existential (traceability) check over region \(R\):}  
We encode the formula
\[
\exists\ \Theta\in R:\ f(\Theta)\le \varepsilon
\]
as an SMT/decision problem. A canonical encoding (assuming \(U\) and \(\mathcal B\) are expressible in the logic) is:
\[
C_R(\Theta)\;\wedge\;
\operatorname{EncodeDistance}(U(\Theta,z),\mathcal B) \le \varepsilon,
\]
where \(\operatorname{EncodeDistance}(\cdot,\cdot)\) is a formula representing the (approximate) distance of \(U(\Theta,z)\) into \(\mathcal B\), or equivalently membership up to tolerance.  
A satisfying assignment \(\Theta^*\) is a witness; it certifies \(z\) is traceable via \(\Theta^*\).

\paragraph{Universal (non-traceability) exclusion check over region \(R\):}  
We encode the negation as
\[
\forall\ \Theta\in R:\ f(\Theta) > \varepsilon,
\]
which is treated in SMT by proving unsatisfiability of the negated existential:
\[
\neg\bigl(\exists\ \Theta\in R:\ f(\Theta)\le \varepsilon\bigr),
\]
i.e., the same formula as above is refuted. In practice, for robustness and efficiency, we strengthen by deriving a \emph{lower bound} \(L_R\) on \(f\) over \(R\) using modulus-of-continuity and samples and then encode the proof obligation:
\[
C_R(\Theta)\;\wedge\;
f(\Theta) < L_R \quad \text{is unsatisfiable, with } L_R > \varepsilon.
\]
Alternatively, one can encode a symbolic certificate (e.g., via interval arithmetic or invariants) asserting \(\min_{\Theta\in R}f(\Theta)\ge L_R>\varepsilon\).

\paragraph{Implementation notes:}
\begin{itemize}[itemsep=1pt]
  \item If \(U\) and \(\mathcal B\) are given by rational arithmetic circuits (EM1), their composition and distance computations can be encoded directly in SMT over reals with rational approximation.  
  \item The modulus of continuity from (EM2) supplies bounds to convert pointwise samples into sound lower/upper bounds over \(R\); these are embedded as side constraints to support universal proofs.  
  \item Use incremental solving: cache previous region proofs to avoid redundant work in refined subdivisions.  
  \item If the solver returns \texttt{sat} for the existential check, extract \(\Theta^*\); if it proves the existential unsatisfiable (or verifies the lower bound certificate), mark \(R\) excluded.  
\end{itemize}


\subsection{Adaptive Traceability Resolution Algorithm}

\begin{remark}[Refinement Fairness]
\label{rem:refinement-fairness}
The algorithm assumes a \emph{fair refinement scheduler}: at each iteration every active region that is not yet excluded is either resolved or subdivided further (e.g., a standard ``refine-all-unresolved'' policy). The existence of such a refinement is proved in Refinement Extension of Guarded Decompositions Theorem (Thm.-\ref{thm:refinement-extension}) in Chapter 2.

Combined with the interior stability margins from Lemma~\ref{lem:branch-stability-margin} (Chapter 2), this guarantees that any region containing a true traceability witness will eventually be refined to diameter \(\le d_{\mathrm{target}}\) and exposed. Note that this spatial refinement fairness is logically distinct from axiom (I) (TI/WF), which governs progression along execution trajectories; the current remark controls region-level search, not the dynamic choice of inputs.
\end{remark}


\begin{algorithm}[H]
\caption{Adaptive Traceability Resolution for input \(z\) with gap \(\gamma>0\)}
\label{alg:adaptive-trace}
\begin{algorithmic}[1]
\Require Input \(z\), certification threshold \(\varepsilon>0\), known gap \(\gamma>0\) such that \(f^*\le \varepsilon\) or \(f^*\ge \varepsilon+\gamma\), initial coarse cover \(\mathcal R_0\) of \(\mathcal M\), basin \(\mathcal B\), modulus of continuity \(\omega(\cdot)\) from Lemma~\ref{lem:modulus-bound}
\Ensure Either: \texttt{TRACEABLE} with witness \(\Theta^*\), or \texttt{NON-TRACEABLE} with gap certificate
\State Compute target refinement diameter \(d_{\mathrm{target}}\) such that \(\omega(d_{\mathrm{target}})<\gamma/2\)
\State \(\mathcal R \gets \mathcal R_0\)
\While{true}
  \ForEach{region \(R\in \mathcal R\)}
    \State Pick representative \(\Theta_0\in R\); compute \(f(\Theta_0)=\operatorname{dist}(U(\Theta_0,z),\mathcal B)\) to required precision via (EM2)/(EM3)
    \If{\(f(\Theta_0) \le \varepsilon + \frac{\gamma}{2}\)} 
      \State \textbf{(Existential success)}: run existential region check (SMT or direct witness extraction) to produce \(\Theta^*\in R\) with \(f(\Theta^*)\le \varepsilon + \frac{\gamma}{2}\)
      \State \Return \texttt{TRACEABLE}, witness \(\Theta^*\)
    \EndIf
    \State Compute lower bound over \(R\): \(L_R := f(\Theta_0) - \omega(\operatorname{diam}(R))\)
    \If{\(L_R > \varepsilon +\frac{\gamma}{2}\)}
      \State \textbf{(Universal exclusion)}: mark \(R\) as excluded
    \EndIf
  \EndFor
  \If{all regions in \(\mathcal R\) are excluded}
    \State \Return \texttt{NON-TRACEABLE} with certificate \(\min_{\Theta} f(\Theta) \ge \varepsilon + \frac{\gamma}{2}\)
  \EndIf
  \State \textbf{Refine:} Replace each non-excluded region \(R\) with a subdivision of smaller-diameter subregions so that eventually \(\operatorname{diam}(R)\le d_{\mathrm{target}}\)
  \State Update \(\mathcal R\) to the new cover
\EndWhile
\end{algorithmic}
\end{algorithm}

\textbf{Notes:}
\begin{itemize}[itemsep=1pt]
  \item The existential check is lightweight when \(f(\Theta_0)\) already certifiably lies within the acceptance band \(\varepsilon + \gamma/2\); a localized search in \(R\) suffices to produce \(\Theta^*\).  
  \item Universal exclusion uses the modulus-of-continuity to lift a sample value to a guaranteed lower bound over the whole region.  
    \item The guarded program provides the logical composition of these region checks: successful existential proofs propagate as accepted branches; excluded regions prune whole subtrees.  
  \item Refinement continues deterministically until either a traceability witness is extracted or every region is excluded; \(\omega(d_{\mathrm{target}})<\gamma/2\) ensures the dichotomy resolves definitively.  
  \item All arithmetic, distance, and membership evaluations are carried out with precision controlled according to (EM2) and (EM3), making the bounds sound.  
  \item Termination and completeness of ATRT rely only on compactness, the gap~$\gamma$, the modulus~$\omega$, and refinement fairness (Remark~\ref{rem:refinement-fairness}); no temporal-input fairness (Axiom (I), TI/WF) is assumed or required.

\end{itemize}


\subsection{Continuity-based Lower Bound Machinery and Invariants}

Under (C) and (EM2), \(U\) is continuous and effectively computable with a known modulus of continuity. Define
\[
f(\Theta) := \operatorname{dist}(U(\Theta,z), \mathcal B).
\]
Since \(U\) is continuous on compact \(\mathcal M\) and \(\mathcal B\) is closed (as a basin shell limit), \(f\) is continuous and attains its minimum.

\begin{lemma}[Effective Local Variation Bound]
\label{lem:modulus-bound}
Suppose (EM2) supplies a polynomial modulus function \(m(k)\) for computing \(U\) to \(k\) bits. Then there exists a computable function \(\omega: \mathbb R_{+}\to\mathbb R_{+}\) (derivable from \(m\) and the Lipschitz behavior induced by finite precision) such that for any \(\Theta,\Theta'\in\mathcal M\),
\[
|f(\Theta)-f(\Theta')| \le \omega\bigl(\operatorname{dist}(\Theta,\Theta')\bigr).
\]
Furthermore, \(\omega\) can be chosen to satisfy \(\omega(r)\to 0\) as \(r\to 0\), and \(\omega\) is computable in time \(\mathrm{poly}(n,\log(1/r))\).
\end{lemma}

\begin{proof}
By (EM2a), each primitive map $F_j$ (either an update $U_j$ or an emit $E_j$) comes with an explicit, effectively computable modulus of continuity $\omega_j:\mathbb{Q}_+\to\mathbb{Q}_+$ satisfying $\omega_j(0)=0$ and
\[
\|x-y\| \le r \quad\Rightarrow\quad \|F_j(x)-F_j(y)\| \le \omega_j(r).
\]
Since $f(\Theta)=\operatorname{dist}(U(\Theta,z),\mathcal{B})$ is obtained by composing finitely many such primitives with each other and with the distance-to-set map, the standard composition rules for moduli of continuity yield an explicit modulus $\omega$ for $f$:
for all $\Theta,\Theta'\in\mathcal{M}$,
\[
\operatorname{dist}(\Theta,\Theta') \le r \quad\Rightarrow\quad |f(\Theta)-f(\Theta')| \le \omega(r).
\]

To make this explicit, note that $U(\cdot,z)$ is the composition of primitive updates along a fixed finite path determined by~$z$, so if $\operatorname{dist}(\Theta,\Theta')\le r$ then
\[
\|U(\Theta,z)-U(\Theta',z)\| \le \omega_U(r)
\]
where $\omega_U$ is obtained by chaining the $\omega_j$ of the primitives in the path.  
The distance function $\operatorname{dist}(\cdot,\mathcal{B})$ is $1$-Lipschitz, hence
\[
|f(\Theta)-f(\Theta')| = \big|\operatorname{dist}(U(\Theta,z),\mathcal{B}) - \operatorname{dist}(U(\Theta',z),\mathcal{B})\big| 
\le \|U(\Theta,z)-U(\Theta',z)\| \le \omega_U(r).
\]
Thus we may take $\omega := \omega_U$ as the desired effective modulus for $f$.

Finally, by (EM2b), evaluating each primitive $F_j$ to $k$ bits of precision costs at most $\mathrm{poly}(k+n)$ steps, where $n$ encodes the instance size. Since the composition depth is fixed, $f(\Theta)$ can be evaluated to any accuracy prescribed by $\omega$ within polynomial time in $k+n$. This completes the proof.
\end{proof}


\begin{corollary}[Region Lower Bound]
\label{cor:region-lower}
Let \(R\subseteq\mathcal M\) be a region of diameter \(\operatorname{diam}(R)\), and let \(\Theta_0\in R\) be a representative sample with computed value \(f(\Theta_0)\). Then for all \(\Theta\in R\),
\[
f(\Theta) \ge f(\Theta_0) - \omega\!\left(\operatorname{diam}(R)\right).
\]
Thus if \(f(\Theta_0) - \omega(\operatorname{diam}(R)) > \varepsilon\), we can safely exclude \(R\) for the traceability check at precision \(\varepsilon\).
\end{corollary}

\section{Appendix to Chapter 5: Deeper analysis of selected theories}

\subsection*{How to read the model audits (template)}
For each theory we check, in order:
\begin{enumerate}[leftmargin=1.25em,itemsep=2pt]
  \item \textbf{In-scope check (ABET availability).} On the laboratory slice $Z^{\mathrm{in}}$ with finite energy/localisation bounds: do (F) Invariant Basin, (S) Certification Standard, and (C) Compact Continuity hold? If not, \emph{Category~1} (FDI out of scope).
  \item \textbf{Effectivity \& structure.} Verify (T) Traceability (open preimages for basin shells), (EM1)–(EM3) (effective primitives), and (I) (no supertasks).
  \item \textbf{Guard construction.} If a \emph{finite} family of guard cells $G_1,\dots,G_m\subseteq Z^{\mathrm{in}}$ covers the in-scope slice with per-guard correctness proofs, the fact is FDI-decomposable $\Rightarrow$ \emph{Category~3}; else \emph{Category~2}.
\end{enumerate}
We also flag a \emph{Traceability Caveat}: if a theory’s latent variables (e.g., a micro-index) are \emph{not} operational inputs, (T) fails and the verdict can drop from Category~3 to Category~2.


%-------------------------------------------------------------
\subsection{Superdeterministic CAI is
            \emph{FDI–decomposable}}
\label{subsec:cai-decomposable}

Let $\Sigma$ be the finite alphabet for a single lattice cell and
$L\gg1$ the (finite) number of cells.  A micro–state is therefore
%
\[
   s \;\in\; \Sigma^{L}, 
   \qquad N:=|\Sigma|^{L}<\infty.
\]
%
The global update $F\colon\Sigma^{L}\!\to\Sigma^{L}$ is \emph{reversible},
so it is a permutation of $\{0,\dots,N-1\}$.  We encode each $s$ by the
binary index $i(s)\in\{0,\dots,N-1\}$ written in
$B=\lceil\log_2 N\rceil=O(L)$ bits.

%..............................................................
\paragraph{Input representation and compact continuity (Gnosis+Structure Axioms).}
The physically accessible slice
$Z^{\mathrm{in}}:=\{0,\dots,N-1\}$ is \emph{finite}, hence compact and
discrete in the product topology used by the FDI framework; guards may
treat micro–indices as ordinary binary strings.

%..............................................................
\paragraph{Deterministic outcome map.}
Fix an apparatus description and time budget $T\!>\!0$.  The pointer
observable is a computable function
\[
    M_T : i(s_0) \;\longmapsto\; a_k \in \{e_1,\dots,e_m\},
\]
obtained by iterating the automaton
$s_t=F^{\,t}(s_0)$ together with the
detector’s finite control logic for $t\le T$.

%..............................................................
\paragraph{Finite guarded decomposition.}
Define \emph{guard sets}
\[
    G_k := M_T^{-1}(e_k) \;\subseteq\; \{0,\dots,N-1\}.
\]
Because $N<\infty$, the family $\{G_k\}_{k=1}^{m}$ is a finite
partition of $Z^{\mathrm{in}}$; the \emph{basin} for outcome $e_k$ is
just $G_k$ itself.  Membership in $G_k$ can be tested in
$O(T+|s_0|)$ time by simulating $F$ for $T$ steps and checking the
pointer, hence satisfies the Polynomial Guard/Basin Test axiom~EM3.

%..............................................................
\paragraph{Certificate and verifier.}
A succinct witness for the tail fact
$\Phi_k\!\!:\,$ “\emph{the pointer shows $e_k$}’’ is the binary string
\[
     w = i(s_0)\quad\text{of length }B=O(L).
\]
Verification algorithm:

\begin{enumerate}[leftmargin=15pt,itemsep=2pt]
 \item Given $(\psi,\text{detector initial macro}, w)$ convert $w$ to
       the micro–state $s_0$ by table lookup.
 \item Simulate $F$ for $T$ steps (reversible CA $\Rightarrow$ each step is local
       and executes in time $O(L)$).
 \item Accept iff the simulated pointer equals $e_k$.
\end{enumerate}

Total runtime is polynomial in $L$; hence
$L_{\Phi_k}\in\mathrm{NP}$ (indeed in
$\mathrm{P}^{F\!}$ if $F$ is provided as an oracle).

%..............................................................
\paragraph{Satisfaction of all FDI axioms.}
\begin{itemize}[leftmargin=12pt,itemsep=1pt]
  \item \textbf{Gnosis+Structure Axioms \& EM1–EM3}: \
        Hold because $Z^{\mathrm{in}}$ is finite and $F$ and $M_T$ are
        explicit, reversible cellular rules.
  \item \textbf{Time Irreversibility (I):} \
        No stochastic branching occurs; simulation halts after $T$
        deterministic steps, the axiom is trivially satisfied.
\end{itemize}

%..............................................................
\paragraph{FDI conclusion.}
Since all axioms are met \emph{and} the guard partition
$\{G_k\}$ is finite, the tail fact $\Phi_k$ is
\emph{decomposable}.  Consequently the certification language sits
inside the ordinary polynomial hierarchy rather than at oracle degree
$0'$.  This places the superdeterministic CA Interpretation firmly in
\textbf{Category~3} as defined in Sec.~\ref{sec:collapse-uncomputable}.
%-------------------------------------------------------------
%=============================================================
\paragraph{Traceability caveat (T).}
The decomposition above treats the micro-index $i(s_0)$ as an \emph{operational input}. If the CAI posits that $i(s_0)$ exists but cannot be accessed (or approximated by an open set of controls) at tick~0, then (T) fails on $Z^{\mathrm{in}}$ and no finite guard family exists. In that case the same dynamics yields \emph{Category~2}, not~3.


\subsection{Certifiability of Convex–Quasi–Linear (CQL) Deterministic
         Collapse}
\label{sec:cql-cert}
%=============================================================

Throughout this section we work in a finite–dimensional truncation
\(H\!\cong\!\mathbb{C}^{d}\) that captures a laboratory’s
finite–energy, finite–particle sector; hence the trace–norm unit ball
\(\mathcal{D}_d\!=\!\{\rho\ge0,\operatorname{Tr}\rho=1\}\) is compact.

%-------------------------------------------------------------
\subsubsection{The Kraus–normalised flow}

Let \(\{\Lambda_t\}_{t\ge0}\) be a \emph{linear}, completely–positive
semigroup
\(\Lambda_t(\rho)=\sum_{\alpha}L_{\alpha}(t)\rho L_{\alpha}^{\dagger}(t)\)
with rational Kraus coefficients, not necessarily trace preserving.
Define the \emph{convex–quasi–linear} flow
\[
   \Phi_t(\rho)
   :=\frac{\Lambda_t(\rho)}{\operatorname{Tr}\Lambda_t(\rho)},
   \qquad 0\le t<\infty .
\]
We assume the \textbf{collapse condition}

\[
   \exists T>0\; \forall\rho\in\mathcal{D}_d:
   \quad \Phi_T(\rho)\in\{\Pi_1,\dots,\Pi_m\},
   \tag{CC}
\]
where \(\{\Pi_k\}\) are one–dimensional projectors onto an orthonormal
pointer basis.  (Examples with this property are constructed in
\cite{RembielinskiCaban2020}.)

The map \(\Phi_T\) is \emph{globally Lipschitz}:
\[
   \|\Phi_T(\rho)-\Phi_T(\sigma)\|_1
   \;\le\; L\,\|\rho-\sigma\|_1,
   \qquad L:=\frac{\|\Lambda_T\|_{1\to1}}{\lambda_{\min}},
\]
where \(\lambda_{\min}:=\min_{\rho}\operatorname{Tr}\Lambda_T(\rho)>0\)
because \(\Lambda_T\) is CP and finite dimensional.

\paragraph{Uniform positivity (well-defined normalisation).}
Assume there exists $c>0$ such that $\operatorname{Tr}\Lambda_T(\rho)\ge c$ for all $\rho\in\mathcal{D}_d$. (This follows, e.g., if $\Lambda_T$ is strictly positive on $\mathcal{D}_d$ or under the collapse condition (CC) together with continuity.) Then $\Phi_T$ is well-defined and globally Lipschitz on $\mathcal{D}_d$:
\[
\|\Phi_T(\rho)-\Phi_T(\sigma)\|_1
\;\le\; \frac{\|\Lambda_T\|_{1\to1}}{c}\;\|\rho-\sigma\|_1
\;=:\; L\;\|\rho-\sigma\|_1.
\]

\paragraph{Finite guarded decomposition (compact cover).}
Each basin $B_k=\Phi_T^{-1}(\{\Pi_k\})$ is closed in the compact $\mathcal{D}_d$, hence compact. For any $0<\varepsilon<\frac{1}{2L}$, the open balls
\[
U_{\rho} \;:=\; \bigl\{\sigma\in \mathcal{D}_d : \|\sigma-\rho\|_1<\varepsilon\bigr\}
\quad (\rho\in B_k)
\]
satisfy $\Phi_T(U_{\rho})\subset B_{1/2}(\Pi_k)$ and therefore decide the pointer as $e_k$. By compactness of $B_k$, extract a finite subcover $U_{\rho_{k,1}},\dots,U_{\rho_{k,r_k}}$. Taking the finite union over $k=1,\dots,m$ yields a finite guard family that covers $\mathcal{D}_d$.


%-------------------------------------------------------------
\subsubsection{Polynomial certificate scheme}

\paragraph{Witness.}
A certificate for the fact
\(\Phi_k:\,“\text{pointer reads }e_k”\) is the triple
\[
  w = \bigl(k,\; \rho_c,\; r\bigr),
\]
where \(k\in\{1,\dots,m\}\), \(\rho_c\in\mathcal{D}_d(\mathbb{Q}[i])\)
has rational entries, and \(r\in\mathbb{Q}_{>0}\) obeys
\(r<\varepsilon_k\) with
\(\|\rho_0-\rho_c\|_1<r\).  The bit–length of \(w\) is
\(O(d^2\log 1/r)\).

\paragraph{Polynomial verification.}
Given a guard label $(k,j)$ and a rational center $\rho_{k,j}$, the verifier checks $\|\rho_0-\rho_{k,j}\|_1<\varepsilon$ and computes $\Phi_T(\rho_{k,j})$ with rational Kraus matrices (EM1), in time polynomial in $d$ and the bit-lengths, then accepts iff $\|\Phi_T(\rho_{k,j})-\Pi_k\|_1<\tfrac{1}{2}$. Hence $L_{\Phi_k}\in\mathbf{NP}$ (and in QCMA with a BQP verifier).

%-------------------------------------------------------------
\subsubsection{Fragility: a single essential discontinuity drops to
            Category 1}

If the flow loses global Lipschitz continuity—e.g.\ there exists
\(\rho^\star\) where \(\Lambda_T(\rho^\star)\) becomes rank-deficient
so that \(\lambda_{\min}=0\)—then (i) guard balls can no longer be made
uniform in radius and (ii) the inverse image of a projector need not be
closed.  Consequently:

\begin{itemize}[leftmargin=12pt,itemsep=1pt]
  \item Representation axiom \textbf{(C)} (compactness under the norm
        topology) fails at \(\rho^\star\).
  \item EM3 is violated because no finite sub–cover can guarantee the
        map to be constant on each guard.
\end{itemize}

The theory therefore exits the FDI domain altogether,
falling to **Category 1** despite being deterministic.
Hence \emph{small analytical pathologies in the CP semigroup can
erase the certifiability advantage}.  This illustrates the delicate
balance required for deterministic collapse models to outperform
Copenhagen in epistemic value.

%-------------------------------------------------------------


%=============================================================
\subsection{Finite-World Many-Interacting-Worlds (MIW)}
\label{subsec:miw-finite}
%=============================================================

We analyse the \emph{finite-$N$} variant of the Hall–Deckert–Wiseman
Many-Interacting-Worlds model, where the universal configuration is a
deterministic swarm of $N\!<\!\infty$ classical “world points’’
$q^{(w)}(t)\in\mathbb R^{3K}$ ($w=1,\dots,N$ for $K$ particles).  The
ordinary Schrödinger predictions are recovered for large but fixed $N$
by suitably choosing the inter-world potential
$V_{N,\gamma}(q^{(1)},\dots,q^{(N)})$ \cite{HallDeckertWiseman2014}.

%-------------------------------------------------------------
\paragraph{State representation and compact input slice (Gnosis+Structure Axioms).}
Assume each particle’s coordinate domain is a bounded box
$[-L,L]^3$ and kinetic energy per particle is $\le E_{\max}$.  Then the
phase-space cube
\[
   Z^{\mathrm{in}}
   :=\Bigl([-L,L]^3\times[-P,P]^3\Bigr)^{\,K N},
   \qquad P:=\sqrt{2mE_{\max}},
\]
is compact and has a straightforward rational-ball naming, satisfying
\textbf{Gnosis+Structure Axioms}.  A micro-state is the tuple
\((q^{(1)},p^{(1)},\dots,q^{(N)},p^{(N)})\) with rational coordinates.

%-------------------------------------------------------------
\paragraph{Deterministic dynamics.}
Time evolution follows a Lipschitz ODE system
\[
   m\ddot q^{(w)} = -\partial_{q^{(w)}}\!
   \bigl[\,V_{\mathrm{ext}}(q^{(w)})
          +V_{N,\gamma}(q^{(1)},\dots,q^{(N)})\bigr]
   \quad (w=1,\dots,N),
\]
whose right-hand side is polynomial if $V_{\mathrm{ext}}$ and
$V_{N,\gamma}$ use rational parameters.  Thus a rational-time step
$\Delta t$ can be simulated with polynomial precision and time, meeting
\textbf{EM1} (effectively computable flow).

%-------------------------------------------------------------
\paragraph{Pointer observable and world-index guard sets.}
Suppose the apparatus records a discrete pointer value
$e_k$ when a configuration $q^{(w)}(T)$ falls inside a rational
detector cell $D_k\subset[-L,L]^{3K}$.  Because trajectories do not
cross, each initial world index $w$ produces \emph{exactly one} pointer
value.  Define

\[
   G_k := \bigl\{\,w\in\{1,\dots,N\}\,:\,
               q^{(w)}(T)\in D_k \bigr\}.
\]

The family $\{G_k\}_{k=1}^{m}$ is a \emph{finite} partition of the
index set, and membership $w\in G_k$ is decidable in
$\mathrm{poly}(K,N)$ time by integrating the ODE for world $w$ alone
(\textbf{EM3}).

%-------------------------------------------------------------
\paragraph{Certificate and verifier.}
A witness for the fact ``pointer reads $e_k$’’ is

\[
   w \;=\; \bigl(k,\; \hat w \in G_k,\;
               (q^{(\hat w)}_0,p^{(\hat w)}_0)\bigr),
\]

where $(q^{(\hat w)}_0,p^{(\hat w)}_0)$ are the rational initial
coordinates of world $\hat w$.  Verification algorithm:

\begin{enumerate}[leftmargin=15pt,itemsep=2pt]
  \item Re-integrate the ODE for world $\hat w$ up to time $T$ with
        rational arithmetic.
  \item Accept iff the resulting $q^{(\hat w)}(T)\in D_k$.
\end{enumerate}

Both steps run in time $\mathrm{poly}(K,N,\log 1/\epsilon)$, so
$L_{\Phi_k}\in\mathrm{NP}$ and, by the standard Kitaev embedding,
$L_{\Phi_k}\in\mathrm{QCMA}$ when a BQP verifier is allowed.

%-------------------------------------------------------------
\paragraph{FDI verdict.}
All Gnosis+Structure Axioms, EM1–EM3, and (I) are satisfied and a
\emph{finite guard} family exists.  Hence every single-outcome fact is
\textbf{FDI-decomposable}; the theory belongs to
\emph{Category 3}—epistemically stronger than Copenhagen.

\paragraph{Traceability caveat (T).}
The guard sets above use the world index $w$ as an operational selector. If the finite-$N$ MIW treats $w$ as a hidden typicality parameter unavailable to the agent at tick~0, then (T) fails (no open preimage in the agent’s input topology), and the theory drops to \emph{Category~2} despite determinism.


%-------------------------------------------------------------
\paragraph{Continuum-world limit drops to Category 2.}
If one lets $N\!\to\!\infty$ so that world indices form a continuum,
the guard sets $G_k$ become uncountable; no finite decomposition is
possible and the model reverts to \textbf{Category 2}.  Thus
\emph{finiteness of $N$ is crucial for certifiability}.
%=============================================================

\end{document}